
\subsubsection{3.1 Neural Functional Networks (NFN)}

We adopt the NFN architecture from Zhou et al. The core building block is NPLinear, which linearly transforms weight tensors while maintaining equivariance. HNPPool produces permutation-invariant representations via pooling over neuron dimensions.

Our NFNRegressor architecture:
\begin{verbatim}
Input weights → NPLinear(32) → ReLU → NPLinear(32) → ReLU → HNPPool → Linear(1)
\end{verbatim}

The model has approximately 70,000--270,000 parameters depending on input dimensionality.

\subsubsection{3.2 MLP-Matched Baseline}

For comparison, we design an MLP that receives flattened weight vectors:
\begin{verbatim}
Flattened weights → Linear(512) → ReLU → Linear(256) → ReLU → Linear(1)
\end{verbatim}

This architecture has substantially more parameters than NFN due to the flattened input dimensionality (approximately 1.8--52 million parameters for the inputs used).

\subsubsection{3.3 Permutation and Invariance Testing}

To verify whether models develop permutation-invariant representations:

\begin{enumerate}
\item Take a trained model M and a test weight tensor W
\item Generate K=10 random permutations $\pi _1, \pi _2$, ..., $\pi$ₖ applied to hidden layer neurons
\item Apply M to each permuted version: $ŷ_i$ = M($\pi _{i}(W))$
\item Compute invariance metrics: mean correlation across predictions, coefficient of variation, maximum deviation
\end{enumerate}

For mathematically invariant models, all predictions should be identical up to floating-point precision.

